\section{All indicators and phase-wise induction masses at multiplicity eight}
\label{sec:1660-all-indicators-phase-masses}

The two saturated modular summands in
Proposition~\ref{prop:1660-forced-third-indicator} determine the complete
indicator sequences, including even orders. Fourier inversion then fixes
the dimension-weighted restriction mass at each central twist. These are
necessary conditions on a hypothetical category; the ring remains
unresolved.

Write $r=\sqrt2$, $d=3+2\sqrt2+2\sqrt3+\sqrt6$, and
$D=4(1+d^2)$. We work with the positive spherical Galois representative
and source order $(1,g,X,gX,\sigma,T)$, where $T=\sigma X$.
Let $\zeta=e^{2\pi i/48}$. The two saturated choices are
\[
 (W_3,W_6)=((33,4),(54,4)),\qquad
 ((36,5),(61,5)),
\]
in the local character-table order of
\S\ref{sec:1660-forced-third-indicator}. They are complex conjugates.
For each choice put
\[
 a_n=2\sqrt3\operatorname{Tr}_{W_3}(P_3s^{-1}t^ns),\qquad
 b_n=2\sqrt6\operatorname{Tr}_{W_6}(P_6s^{-1}t^ns),
\]
where $P_j$ projects onto the one-dimensional $t$-fixed line of $W_j$.

\begin{theorem}[Complete indicator sequences]
\label{thm:1660-all-indicators}
For every positive integer $n$, put
\[
 g_n=\begin{cases}1&2\mid n,\\0&2\nmid n,\end{cases}
 \qquad
 s_n=\begin{cases}\sqrt2\cos(3n\pi/8)&2\mid n,\\0&2\nmid n.\end{cases}
\]
Then
\begin{align*}
 \nu_n(X)=\nu_n(gX)&=\frac{3+3g_n+4s_n+2a_n+b_n}{2},\\
 \nu_n(T)&=2+2g_n+3s_n+a_n+b_n.
\end{align*}
Together with $\nu_n(1)=1$, $\nu_n(g)=g_n$ and
$\nu_n(\sigma)=s_n$, these give all six indicator sequences.
Every sequence is periodic of period dividing $48$.
\end{theorem}
\begin{proof}
The projection and central-indicator identities in the proof of
Proposition~\ref{prop:1660-forced-third-indicator} hold for every $n$;
only its subsequent simplification used oddness. In particular,
\[
 \sum_i\ell_j(\dim X_i)\nu_n(X_i)
 =\frac{\operatorname{Tr}_{W_j}(P_js^{-1}t^ns)}{4\sqrt j}
 \qquad(j=3,6).
\]
Saturation makes the orthogonal projection of $e_1$ onto $W_j$ equal to
$w_j/4$, so no unknown complementary summand contributes to this identity.
The tensor involution of Proposition~\ref{prop:1660-tensor-involution}
gives equal indicators for $X$ and $gX$ at every order.
The invertible $g$ has trivial pointed associator and indicator $+1$
at even orders. For the positive Ising subcategory with
$\nu_2(\sigma)=-1$, the modular indicator formula
\cite[Theorem~7.5]{NS} gives $s_n$: using an Ising braiding with
$\theta_\sigma=e^{3\pi i/8}$, the four fusion pairs contributing to
$\sigma$ yield zero at odd orders and
$\sqrt2\cos(3n\pi/8)$ at even orders. This calculation depends only on
the underlying positive pivotal Ising category.

Substitute these indicators into the two coefficient rows
\[
 24\ell_3=(-1,-1,1,1,-1,-1),\qquad
 48\ell_6=(-1,-1,-1,-1,-2,2).
\]
Writing $x=\nu_n(X)=\nu_n(gX)$ and $u=\nu_n(T)$ gives
\[
 2x-u=a_n+1+g_n+s_n,\qquad
 -2x+2u=b_n+1+g_n+2s_n.
\]
Solving gives the stated formulas. Every actual central twist has order
dividing $48$ by Proposition~\ref{prop:1660-uniform-indicators}; the
central-indicator formula therefore proves the periodicity as well.
\end{proof}

\begin{theorem}[Central phase masses]
\label{thm:1660-phase-masses}
For $k\in\mathbb Z/48$ and a base simple $X_i$, define
\[
 M_i(k)=\sum_{\theta_Z=\zeta^k}\dim(Z)[F(Z):X_i],
 \qquad Z\in\operatorname{Irr}Z(\mathcal C).
\]
These masses are determined by
\begin{equation}\label{eq:1660-phase-fourier}
 M_i(k)=\frac D{48}\sum_{n=1}^{48}\nu_n(X_i)\zeta^{-kn}.
\end{equation}
For the choice $((33,4),(54,4))$, the complete twist support is
\[
 \{0,8,9,12,14,15,20,24,26,30,32,33,36,39,42,44\}.
\]
For the conjugate choice, negate these exponents modulo $48$.
There are exactly two simples in each of the four phases
$9,15,33,39$. All eight have the same restriction
\begin{equation}\label{eq:1660-eight-spin-columns}
 F(Z)=2X+2gX+\sigma+3T.
\end{equation}
These are precisely the central simples whose restrictions contain
$\sigma$.
\end{theorem}
\begin{proof}
The center formula \cite[Theorem~4.1]{NS} and adjunction give
\[
 D\nu_n(X_i)=\sum_k M_i(k)\zeta^{kn}.
\]
Discrete Fourier inversion gives \eqref{eq:1660-phase-fourier}.
The exact character averages determining $a_n,b_n$ are
\[
 \frac1m\sum_{j=0}^{m-1}
 \operatorname{char}_{W}\!\begin{pmatrix}1-jn&j\\-n&1\end{pmatrix}
 \quad(m=16,3),
\]
multiplied across the two local factors. Evaluation of these finite
sums and the Fourier transform in $\mathbb Q(\zeta)$ gives the displayed
support. Every mass outside it is zero. Positive dimensions and
faithfulness of $F$ imply that no central simple can lie at a phase
with all six masses zero.

For each $k=9,15,33,39$, the nonzero coordinates relevant here are
\[
 \frac{(M_X(k),M_\sigma(k),M_T(k))}{d}
 =\bigl(8\sqrt3+8\sqrt6,\;4\sqrt3+4\sqrt6,\;
        12\sqrt3+12\sqrt6\bigr).
\]
The $gX$ coordinate equals the $X$ coordinate, and every other phase
has $M_\sigma(k)=0$. Proposition~\ref{prop:1660-uniform-indicators}
shows that a $\sigma$-containing column is
$(0,0,b,c,1,3)$ with $b+c=4$, and its dimension divided by $d$ is
$2\sqrt3+2\sqrt6$. Hence exactly two such simples occur in each
of the four phases.

Such a simple cannot be self-dual. A self-duality form on its
restriction would have to have sign $-1$ because the multiplicity
of $\sigma$ is odd, and sign $+1$ because the multiplicity of $T$
is odd. Duality preserves both twists and restrictions, so the two
$\sigma$-containing objects in a phase are a dual pair with identical
columns. Their fixed total $b+c=4$ exhausts the entire displayed
$M_X+M_{gX}$ mass. Positivity therefore leaves no $X$ or $gX$
contribution from any other object in that phase. The separate $X$
and $gX$ masses now give $b=c=2$. These two columns also exhaust the
$\sigma$ and $T$ masses, while the unit and $g$ masses are zero.
Faithfulness and positivity exclude every additional simple in the
phase, proving \eqref{eq:1660-eight-spin-columns} and the exact fiber size.
\end{proof}

\begin{corollary}[Homogeneous phase Gram constraints]
\label{cor:1660-homogeneous-phase-grams}
Use the first conjugate choice in Theorem~\ref{thm:1660-phase-masses}.
At any phase in the following table, write the restriction columns as
$(0,0,b,c,0,a)$ and put $q=b+c$. The phase-wise sums satisfy
\[
\begin{array}{c|rrr}
 k&\sum q^2&\sum qa&\sum a^2\\\hline
 8,32&40&24&16\\
 12&4&4&4\\
 14,26&16&16&16\\
 20&72&40&24\\
 30,42,44&8&8&8\\
 36&36&20&12
\end{array}
\]
Moreover $\sum bq=\sum cq$ and $\sum ba=\sum ca$ in every
listed phase. Every individual column in phases
$12,14,26,30,42,44$ has $q=a$.
\end{corollary}
\begin{proof}
The unit can occur in a restriction only at twist one. Since
$I(g)=hI(1)$ for the central fermion $h$, its constituents have twists
$1$ or $-1$; thus the listed phases have no $g$ coordinate either.
Theorem~\ref{thm:1660-phase-masses} removes the $\sigma$ coordinate.
Each remaining column has dimension $d(q+\sqrt2a)$. Consequently
\[
 M_X/d=\sum bq+\sqrt2\sum ba,\quad
 M_{gX}/d=\sum cq+\sqrt2\sum ca,\quad
 M_T/d=\sum aq+\sqrt2\sum a^2.
\]
The exact masses have $M_X=M_{gX}$. Comparing the rational and
$\sqrt2$ coefficients gives the equalities and the table. In its
homogeneous rows, $\sum(q-a)^2=0$. Every summand is a nonnegative
integer, so each is zero.
\end{proof}

\paragraph{Exact verification and scope.}
The source-bound evidence in
\filename{evidence/rigidity\_20261017/rank6/r1660/} regenerates the
local character averages, all $48$ indicator values, and all $288$
phase masses for each conjugate choice. An independent SymPy computation
uses polynomial remainders modulo $\Phi_{48}$ and rational linear solves;
it checks every Fourier inversion, every inverse transform, the full
Ising indicator sequence, and independence of the dimension-field basis.
The checks certify the finite arithmetic in these arguments, not
categorical existence or nonexistence. In particular, no unknown modular
complement is assumed to coincide with the rank-$37$ witness.

\begin{corollary}[The full phase-mass restriction domain]
\label{cor:1660-phase-640}
The $1136$ labelled restriction profiles satisfying the preceding
Galois--fermion and graded-induction constraints have exactly $640$
profiles compatible with all the phase Gram identities above. Their
three branch sizes are $68,286,286$.
\end{corollary}
\begin{proof}
In each low frame, the six unit-induced columns exhaust all six
phase-zero masses in \eqref{eq:1660-phase-fourier}. Their four forced
even fermion partners exhaust all six phase-$24$ masses. These are
exact equalities in $\mathbb Q(\sqrt2,\sqrt3)$; positivity excludes
any additional column at either phase. Remove these four high partners.
Theorem~\ref{thm:1660-phase-masses} requires exactly eight
spin columns $(0,2,2,1,3)$ in the high-column convention that omits the
unit coordinate. For each remaining column type $(b,c,a)$, its five
phase-moment contributions are
\[
 \bigl(b(b+c),\ c(b+c),\ ba,\ ca,\ a^2\bigr).
\]
Allocate the available multiplicities among the ten phases in
Corollary~\ref{cor:1660-homogeneous-phase-grams}, with target vector
$(Q/2,Q/2,H/2,H/2,A)$ when the displayed triple is $(Q,H,A)$.
All variables are nonnegative integers bounded by the available
column multiplicities. The retained witnesses give every passing
allocation. An independent exact coefficient-convolution calculation
constructs all possible phase monomials; meet-in-the-middle support
products then test whether their product has the required global
multiplicity vector. It reproduces precisely the three stated counts.
Both implementations use exact integers and exhaust their finite domains.
This tests necessary restriction data; it imposes no existence of
associators or full modular fusion matrices.
\end{proof}

